\section{Discussion and Conclusion}
\label{sec:discussion}

\subsection{What generalizes and what does not}
\label{sec:discussion-general}

Some of our results follow from the mathematics of redundancy and should hold
for any knowledge graph and any path-based task.
\begin{itemize}[leftmargin=*,itemsep=2pt]
\item Under blind damage, the survival of a bounded-hop query is a reliability
polynomial in the retention probability, Eq.~\eqref{eq:r2}. It is set by the
query's independent routes and has no threshold. A giant component that is far
from its threshold therefore says little about such queries.
\item Connectivity, reachability and bounded-hop retrieval are instances of one
expression, Eq.~\eqref{eq:principle}. They differ in the unit and in the low end
of the redundancy distribution. The quantity with the most single-unit objects
fails first.
\item Damage aimed at hubs removes the redundant units that sit on hubs.
Counting units then overstates protection, and weighting each unit by its
degree restores the prediction. The attack penalty measures how much a task
depends on hubs.
\item The best repair depends on the horizon. Layer-aware repair is best for
short typed paths at large budgets, and structural repair is best when longer
paths are allowed or when the budget is small.
\end{itemize}

Other results are specific to our sample and benchmark.
\begin{itemize}[leftmargin=*,itemsep=2pt]
\item The magnitudes: 98\% connectivity at 50\% edge loss and the 57-point gap.
The full filtered graph is less robust under random damage. For all of its
connected drug--disease pairs, Eq.~\eqref{eq:r2} gives $R_2=0.30$ at $f=0.5$,
rather than the benchmark's 0.41. Relative to theory, its excess fragility
under attack is similar to that of the sample.
\item The drug-side bottleneck at long horizons. It depends on how many drugs
have a single target, which is 69\% in the sample and 53\% in the full graph.
\item The hub-heavy routes, which follow partly from how the benchmark was
built: a shared protein must link both a drug and a disease.
\end{itemize}

\subsection{Limitations}
\label{sec:discussion-limits}

\begin{itemize}[leftmargin=*,itemsep=2pt]
\item \textbf{Sample.} The standard graph is a BFS sample around one disease.
It is 93\% proteins and contains 74 of the 6,282 drugs. We checked structural
claims on the full filtered graph. Simulated retrieval exists only on the sample,
although for $L=2$ the exact formula extends to all pairs of the full graph.
\item \textbf{Benchmark.} A query succeeds if a typed path exists. This is a
self-referential criterion with no external ground truth: PrimeKG's
drug--indication edges were not used. The queries were also selected to
succeed initially and to balance redundancy.
\item \textbf{Damage model.} Independent random deletion is a clean null model,
not a realistic failure process. Real knowledge graphs lose or miss facts
through study bias and curation choices, which are correlated with degree. Each
attack is a single deterministic order computed once on the intact graph.
Betweenness is estimated from 100 sampled sources.
\item \textbf{Theory.} Equations~\eqref{eq:gf-bond}, \eqref{eq:rinf}
and~\eqref{eq:mincut} assume a locally tree-like, uncorrelated graph. The null
models identify degree correlations as one cause of failure under attack, but
they do not resolve the remaining half. The Molloy--Reed threshold is an
estimate that our damage grid does not reach.
\item \textbf{Repair.} The task-aware score is static and saturates. Only 60
test queries are evaluated, which behave like about 27 independent ones, and 19
of the 21 test drugs also appear in training queries. Of 24 paired comparisons,
only a few survive a correction for multiple testing.
\end{itemize}

\subsection{Conclusion}
\label{sec:discussion-conclusion}

A knowledge graph can stay almost fully connected while losing most of the
short paths that its queries use. On a PrimeKG sample, half of the edges can be
removed with 98\% of nodes still connected, while only 41\% of two-hop queries
survive. The reason is that each quantity is protected by its own redundant
units. Nodes have many independent edges, but many queries have a single
shared protein. One expression, one minus the generating function of the
redundancy distribution, describes all three observables. It is exact for
two-hop retrieval under blind damage and approximate but accurate for
long-range reachability. Under targeted damage, counting units is no longer
enough, because the units sit on hubs. Weighting each unit by its degree
restores the prediction. For anyone who maintains or evaluates a
knowledge graph, the practical advice is to measure robustness on the task
itself, to count independent routes rather than paths, to ask how many of those
routes pass through hubs, and to report the graph, queries, horizon and
damage model together.
